% Preamble of the INDI report. Build with lualatex (see build.sh).
\documentclass[11pt,a4paper,twoside=false,parskip=false,headings=normal,numbers=noenddot]{scrreprt}

% ---------------------------------------------------------------- fonts
\usepackage{fontspec}
\IfFontExistsTF{Libertinus Serif}{%
  \setmainfont{Libertinus Serif}%
  \setsansfont{Libertinus Sans}%
  \IfFontExistsTF{Libertinus Mono}{\setmonofont[Scale=MatchLowercase]{Libertinus Mono}}%
    {\setmonofont[Scale=MatchLowercase]{Latin Modern Mono}}%
  \usepackage{unicode-math}%
  \setmathfont{Libertinus Math}%
  \setmathfont[version=bold]{Libertinus Math}%
}{%
  \setmainfont{Latin Modern Roman}%
  \setsansfont{Latin Modern Sans}%
  \setmonofont{Latin Modern Mono}%
  \usepackage{unicode-math}%
  \setmathfont{Latin Modern Math}%
}
\usepackage{microtype}
\usepackage[english]{babel}

% ---------------------------------------------------------------- maths, tables, graphics
\usepackage{amsmath}  % unicode-math supplies the symbols of amssymb
\usepackage{booktabs,tabularx,array}
\usepackage{graphicx}
\usepackage[export]{adjustbox}
\usepackage{xcolor}
\usepackage{enumitem}
\setlist{nosep,leftmargin=1.4em}
\usepackage{tikz}
\usetikzlibrary{positioning,arrows.meta,calc,fit,backgrounds,shapes.geometric}
\usepackage{pgfplots}
\pgfplotsset{compat=1.17}
\usepackage[labelfont={bf,color=indiblue},font=small,format=hang,labelsep=colon]{caption}
\usepackage{subcaption}

% ---------------------------------------------------------------- palette
\definecolor{indiblue}{HTML}{1F4E79}
\definecolor{indiorange}{HTML}{C55A11}
\definecolor{indigrey}{HTML}{595959}
\definecolor{indilight}{HTML}{DCE6F2}
\definecolor{indigreen}{HTML}{548235}
\definecolor{indired}{HTML}{A50021}

% ---------------------------------------------------------------- headings
\addtokomafont{disposition}{\sffamily\color{indiblue}}
\setkomafont{chapter}{\sffamily\bfseries\color{indiblue}\LARGE}
\setkomafont{section}{\sffamily\bfseries\color{indiblue}\Large}
\setkomafont{subsection}{\sffamily\bfseries\color{indiblue}\large}
\setkomafont{subsubsection}{\sffamily\bfseries\color{indiblue}}
\setkomafont{paragraph}{\sffamily\bfseries\color{indigrey}}
\setkomafont{descriptionlabel}{\sffamily\bfseries\color{indiblue}}
\setkomafont{pageheadfoot}{\sffamily\small\color{indigrey}}
\setkomafont{pagenumber}{\sffamily\small\color{indigrey}}

% ---------------------------------------------------------------- headers
\usepackage[headwidth=\textwidth,footwidth=\textwidth]{scrlayer-scrpage}
\automark[section]{chapter}
\pagestyle{scrheadings}
\clearpairofpagestyles
\ihead{\headmark}
\ohead{\pagemark}
\setkomafont{headsepline}{\color{indiblue}}
\KOMAoption{headsepline}{true}
\renewcommand*{\chaptermarkformat}{\chapapp~\thechapter\autodot\enskip}
\renewcommand*{\sectionmarkformat}{\thesection\autodot\enskip}
\renewcommand*{\chapterpagestyle}{scrplain}
\pagestyle{scrheadings}
\ofoot[\pagemark]{}
\cfoot[]{}

% ---------------------------------------------------------------- links and references
\usepackage[colorlinks=true,linkcolor=indiblue,citecolor=indiblue,urlcolor=indiblue,
  pdftitle={Incremental Nonlinear Dynamic Inversion for PX4 Multicopters},
  pdfsubject={Inner rate loop and outer acceleration loop: design, analysis and simulation validation},
  pdfkeywords={INDI, PX4, multicopter, rate control, acceleration control, simulation},
  pdfauthor={},
  bookmarksnumbered=true,pdfusetitle=false]{hyperref}
\usepackage[capitalise,nameinlink]{cleveref}
\crefname{chapter}{Chapter}{Chapters}
\crefname{appendix}{Appendix}{Appendices}
\crefname{section}{Section}{Sections}
\crefname{figure}{Figure}{Figures}
\crefname{table}{Table}{Tables}
\crefname{equation}{Equation}{Equations}
\creflabelformat{equation}{#2\textup{(#1)}#3}

% ---------------------------------------------------------------- numbering
\numberwithin{equation}{chapter}
\counterwithin{figure}{chapter}
\counterwithin{table}{chapter}

% ---------------------------------------------------------------- boxes
\usepackage{tcolorbox}
\tcbuselibrary{breakable,skins}
\newtcolorbox{keyresult}[1][Key result]{%
  enhanced,breakable,colback=indilight!60,colframe=indiblue,
  coltitle=white,fonttitle=\sffamily\bfseries\small,title={#1},
  boxrule=0.6pt,arc=1.5pt,left=6pt,right=6pt,top=4pt,bottom=4pt,
  before skip=10pt,after skip=10pt}
\newtcolorbox{indinote}[1][Note]{%
  enhanced,breakable,colback=white,colframe=indiorange,
  coltitle=white,fonttitle=\sffamily\bfseries\small,title={#1},
  boxrule=0.6pt,arc=1.5pt,left=6pt,right=6pt,top=4pt,bottom=4pt,
  before skip=10pt,after skip=10pt}

% ---------------------------------------------------------------- macros: parameters and topics
% \param{MC_INDI_EN}: monospace, a line break may follow every underscore.
\ExplSyntaxOn
\tl_new:N \l__indi_tmp_tl
\tl_new:N \l__indi_us_tl
\tl_set:Nx \l__indi_us_tl { \char_generate:nn { 95 } { 8 } }
\NewDocumentCommand{\indiBreakUs}{m}
  {
    \tl_set:Nn \l__indi_tmp_tl {#1}
    \tl_replace_all:NVn \l__indi_tmp_tl \l__indi_us_tl { \textunderscore \hspace{0pt} }
    \l__indi_tmp_tl
  }
\ExplSyntaxOff
\newcommand{\param}[1]{\texttt{\indiBreakUs{#1}}}
\newcommand{\topic}[1]{\texttt{\indiBreakUs{#1}}}
\newcommand{\module}[1]{\texttt{\indiBreakUs{#1}}}

% ---------------------------------------------------------------- macros: diagrams and figures
\IfFileExists{diagrams/diagram_style.tex}{% Shared TikZ styles for the INDI report diagrams.
% Needs tikz with libraries: positioning, arrows.meta, calc, fit, backgrounds, shapes.geometric
\providecolor{indiblue}{HTML}{1F4E79}
\providecolor{indiorange}{HTML}{C55A11}
\providecolor{indigrey}{HTML}{595959}
\providecolor{indilight}{HTML}{DCE6F2}
\providecolor{indigreen}{HTML}{548235}
\providecolor{indired}{HTML}{A50021}

\providecommand{\indimono}[1]{{\ttfamily\scriptsize #1}}

\tikzset{
  % ---- text helpers ----
  dtext/.style   = {font=\sffamily\footnotesize, align=center},
  mono/.style    = {font=\ttfamily\scriptsize},
  lab/.style     = {font=\sffamily\scriptsize, text=indigrey, inner sep=1.2pt},
  labm/.style    = {font=\ttfamily\fontsize{6.5}{7.5}\selectfont, text=indigrey, inner sep=1.2pt},
  sgn/.style     = {font=\sffamily\scriptsize\bfseries, inner sep=0.8pt, text=indiblue},
  % ---- arrows ----
  >=Stealth,
  signal/.style  = {-{Stealth[length=1.9mm,width=1.5mm]}, line width=0.5pt, draw=indiblue},
  measured/.style= {-{Stealth[length=1.9mm,width=1.5mm]}, line width=0.5pt, draw=indiorange},
  modelled/.style= {-{Stealth[length=1.9mm,width=1.5mm]}, line width=0.5pt, draw=indigrey, dash pattern=on 3pt off 2pt},
  plain/.style   = {line width=0.5pt, draw=indiblue},
  % ---- blocks ----
  block/.style   = {draw=indiblue, fill=indilight, line width=0.5pt, rounded corners=2.5pt,
                    minimum height=0.95cm, inner sep=3pt, dtext},
  plant/.style   = {draw=indigrey, fill=indigrey!14, line width=0.5pt, rounded corners=2.5pt,
                    minimum height=0.95cm, inner sep=3pt, dtext},
  measblock/.style = {draw=indiorange, fill=indiorange!10, line width=0.5pt, rounded corners=2.5pt,
                    minimum height=0.95cm, inner sep=3pt, dtext},
  modblock/.style  = {draw=indigrey, fill=white, line width=0.5pt, rounded corners=2.5pt, dash pattern=on 3pt off 2pt,
                    minimum height=0.95cm, inner sep=3pt, dtext},
  sum/.style     = {circle, draw=indiblue, fill=white, line width=0.5pt, inner sep=0pt,
                    minimum size=0.55cm, font=\sffamily\scriptsize},
  switch/.style  = {draw=indiblue, fill=white, line width=0.5pt, rounded corners=2.5pt,
                    minimum height=0.8cm, inner sep=3pt, dtext},
  junction/.style= {circle, fill=indiblue, inner sep=0pt, minimum size=1.5mm},
  junctionm/.style={circle, fill=indiorange, inner sep=0pt, minimum size=1.5mm},
  badge/.style   = {circle, draw=indigreen, fill=white, text=indigreen, line width=0.5pt, inner sep=0.5pt,
                    minimum size=0.42cm, font=\sffamily\scriptsize\bfseries},
  % ---- architecture ----
  module-stock/.style = {draw=indigrey, fill=indigrey!10, line width=0.5pt, rounded corners=2.5pt,
                    inner sep=3pt, align=center, font=\sffamily\footnotesize, minimum height=1.35cm},
  module-indi/.style  = {draw=indiblue, fill=indilight, line width=1.1pt, rounded corners=2.5pt,
                    inner sep=3pt, align=center, font=\sffamily\footnotesize, minimum height=1.35cm},
  module-kept/.style  = {draw=indigrey, fill=indigrey!10, line width=0.5pt, rounded corners=2.5pt,
                    inner sep=3pt, align=center, font=\sffamily\footnotesize, minimum height=1.0cm},
  source/.style  = {draw=indigrey, fill=white, line width=0.5pt, rounded corners=6pt,
                    inner sep=3pt, align=center, font=\sffamily\footnotesize, minimum height=0.85cm},
  group/.style   = {draw=indigrey, line width=0.4pt, dash pattern=on 2pt off 2pt, rounded corners=4pt,
                    inner sep=0pt},
  grouplabel/.style = {font=\sffamily\scriptsize\itshape, text=indigrey, inner sep=2pt},
  % ---- state diagrams ----
  state/.style   = {draw=indiblue, fill=indilight, line width=0.5pt, rounded corners=6pt,
                    minimum height=1.0cm, inner sep=3pt, dtext},
  stateflight/.style = {draw=indigreen, fill=indigreen!12, line width=0.5pt, rounded corners=6pt,
                    minimum height=1.0cm, inner sep=3pt, dtext},
}
}{}

\newcommand{\indiplaceholder}[2]{%
  \fbox{\parbox[c][3.2cm][c]{0.9\linewidth}{\centering\sffamily\color{indigrey}
  Placeholder: #1\\ (#2 not yet generated)}}}

% \indidiagram{name}{caption}{label}: input diagrams/name.tex (a bare tikzpicture) in a figure
\newcommand{\indidiagram}[3]{%
  \begin{figure}[tbp]
    \centering
    \IfFileExists{diagrams/#1.tex}{\adjustbox{max width=\linewidth}{\input{diagrams/#1.tex}}}{\indiplaceholder{#1}{diagrams/#1.tex}}
    \caption{#2}\label{#3}
  \end{figure}}

% \indifigure{name}{caption}{label}: include figures/name.pdf in a figure
\newcommand{\indifigure}[3]{%
  \begin{figure}[tbp]
    \centering
    \IfFileExists{figures/#1.pdf}{\includegraphics[width=\linewidth]{figures/#1.pdf}}{\indiplaceholder{#1}{figures/#1.pdf}}
    \caption{#2}\label{#3}
  \end{figure}}

% ---------------------------------------------------------------- notation macros
\newcommand{\vect}[1]{\symbf{#1}}
\newcommand{\dd}{\mathrm{d}}
\newcommand{\ts}{T_{\mathrm{s}}}
\newcommand{\fc}{f_{\mathrm{c}}}
\newcommand{\aref}{\mathbf{a}_{\mathrm{ref}}}
\newcommand{\Omax}{\Omega_{\max}}

% ---------------------------------------------------------------- table helpers
\renewcommand{\arraystretch}{1.15}
\newcolumntype{L}[1]{>{\raggedright\arraybackslash}p{#1}}
\newcolumntype{Y}{>{\raggedright\arraybackslash}X}
\newcommand{\passed}{\textcolor{indigreen}{\textbf{PASSED}}}
\newcommand{\notpassed}{\textcolor{indired}{\textbf{NOT PASSED}}}

\setlength{\emergencystretch}{2em}
\raggedbottom
